\section{Evaluation}
\label{sec:evaluation}

\subsection{Environment and methodology}

All Linux measurements were performed on an Intel Core i5-14400F with 16 threads,
running a Fedora-class Linux x86-64 host. The stable toolchain is Rust 1.99.0,
pinned by \code{rust-toolchain.toml}; Kani is 0.68.0; Miri and libFuzzer use the
nightly toolchain of the 2026-08 family; Ghidra is 12.2; radare2 is 5.9.8; angr
is 10.0.7; KLEE is 3.3; Creusot is 0.15; GAMBA, SiMBA, and MBA-Blast are the
upstream tools described in Section~\ref{sec:background}.

The evaluation distinguishes four kinds of evidence:

\begin{enumerate}
  \item \textbf{Execution evidence}: a real protected binary produces the same
        observable output as the native binary on the tested inputs.
  \item \textbf{Safety evidence}: a bounded proof, dynamic analysis, or fuzz
        campaign covers a specific component or input class.
  \item \textbf{Recovery-cost evidence}: an external analysis tool is run for a
        fixed budget, and the result, including timeouts and failures, is
        recorded.
  \item \textbf{Diversity evidence}: two builds are compared bytewise and at the
        level of program structure.
\end{enumerate}

No result is extrapolated from a different binary, tool version, or host. The
commands and outputs are recorded under \path{docs/VERIFICATION.md}; the harness
sources are in the repository.

\subsection{Functional correctness}

The workspace runs 226 tests in debug mode and 226 tests in release mode, with
two ignored measurement hooks (MBA corpus emission and statistical timing). The
release-mode run is intentional: the project's history contains a release-only
failure in which an unimported semantic constant became a catch-all match binding
in the generated engine. Release-mode tests and the build-script dependency
tracking now cover that class of bug.

The end-to-end fixtures compile real C or Rust victims, protect them, execute the
protected binary, and compare outputs with native:

\begin{itemize}
  \item a branch/loop/overflow victim;
  \item a RIP-relative read-only table victim;
  \item an eight-argument System V stack-argument victim;
  \item SSE2 move and arithmetic victims;
  \item a VEX-encoded 128-bit AVX integer victim;
  \item a repeated-string \code{movs}/\code{stos} victim;
  \item a byte-swap victim in normal and constant-time modes;
  \item a direct caller/callee pair and a \verb|--follow-calls| pair;
  \item an eight-thread concurrent reentrancy victim;
  \item a SIGALRM handler reentrancy victim;
  \item a constant-time branch/store victim whose inactive branch must not write
        memory.
\end{itemize}

The differential stress harness compiles six functions covering arithmetic,
direct calls, RIP-relative data, AVX/SSE, and repeated strings. Over 10,000
iterations in the fast profile and 1,000 iterations in the max profile, the
protected binary matched the native output exactly. A constant-time build of the
arithmetic, RIP-data, and SSE functions matched native over 1,000 iterations.

The release mode is also tested with the full workspace suite, not only the
generated runtime. This is important because the protector itself is a release
binary in deployment, and its behavior can differ from the debug build.

\subsection{Safety and robustness}

\subsubsection{Bounded proofs}

Kani verifies five core harnesses and four VM harnesses:

\begin{itemize}
  \item core arithmetic and expression helpers;
  \item bounded memory load/store helpers;
  \item masked register round-trip;
  \item masked memory round-trip.
\end{itemize}

The decoder is too large for a useful CBMC run on this host; it is covered by an
opcode/error corpus, property tests, Miri, and libFuzzer. This is a proof-scope
limitation, not a claim that the decoder is proven.

\subsubsection{Miri}

Miri runs the \code{sirus-vm} library tests: 11/11 pass on the release tree. The
full VM execution suite and the core suite are opt-in and slow; they were not
rerun for every release candidate. The library suite checks the Rust abstract
machine for the engine's memory and arithmetic helpers.

\subsubsection{Fuzzing}

The libFuzzer targets exercise the VM payload decoder, the lifter/compiler, and
the ELF/PE parsers. Table~\ref{tab:fuzz} reports the final-tree sessions and the
cumulative totals across recorded sessions.

\begin{table}[htbp]
\centering
\small
\caption{libFuzzer coverage on the release tree and cumulative recorded runs.}
\label{tab:fuzz}
\begin{tabular}{@{}lrr@{}}
\toprule
Target & Final-tree runs & Cumulative runs \\
\midrule
VM payload & 789,455 & 4,883,875 \\
Lifter/compiler & 133,973 & 405,768 \\
ELF parser & 13,516,906 & 33,588,394 \\
PE parser & 26,539,167 & 55,510,941 \\
\bottomrule
\end{tabular}
\end{table}

No crash, panic, out-of-bounds access, or timeout was observed. The VM target
constructs a valid integrity key for arbitrary payloads, so the decoder body path
is exercised rather than only the integrity rejection path. The in-tree
deterministic fuzzer additionally runs 20,000 random payloads and 10,000
mutations of a valid payload. The malformed-payload suite covers bad counts,
offsets, lengths, opcodes, registers, traps, and fuel exhaustion.

Valgrind reports no runtime memory errors on the protected sample. The
\path{scripts/red-team.sh} harness also checks body scrubbing, payload entropy,
string absence, tamper behavior, and forced-return behavior.

\subsection{Static analysis}

\subsubsection{Ghidra}

Ghidra 12.2 headless analyzed the protected red-team binary. It identified 19
scoped functions: 17 in the injected \code{.sirus} segment and two patched host
targets. It completed 18 of 19 decompilations. The two host targets reduce to
\code{halt_baddata()} after the jump/scrub; the injected segment yields VM
engine and state machinery. One internal \code{.sirus} function hit a Ghidra
pcode/unreadable-memory warning and did not complete. This is recorded as a
decompiler limitation, not as a Sirus correctness claim.

The host binary as a whole retains Rust/std strings and symbol names; the scoped
\code{.sirus} string scan is clean. Sirus does not strip the host image, and the
paper does not claim that it does.

\subsubsection{radare2}

radare2 5.9.8 \code{pdc} reduces the two patched host functions to a jump/goto
stub and a return. No original logic appears in the decompiled output.

\subsection{Symbolic execution}

angr 10.0.7 loads the protected ELF and starts at the patched
\code{sirus_demo_check} symbol with two symbolic 64-bit arguments. The requested
target is the reference output for $(7,3)$. Within a 120-second budget the
manager has five active states and no deadended or errored state; no preimage is
found. This is a bounded timeout, not a proof of unsatisfiability. Unicorn
support was unavailable in the Python 3.13 environment, so angr used its VEX
engine.

KLEE 3.3 runs the generated VM source as LLVM 18 bitcode inside the official
container. The harness makes two arguments symbolic and calls the entry. The
final run executed 7,234,878 instructions, completed zero paths, left eight
partially completed paths, and generated five tests. It reported no memory
errors. The run is a bounded symbolic probe; the lack of a completed path is not
a proof that no path exists.

\subsection{MBA simplification}

The external MBA study uses the exact expressions emitted after the production
MBA pass. Thirty cases are generated. Table~\ref{tab:mba} summarizes the result.

\begin{table}[htbp]
\centering
\small
\caption{External MBA simplification results on the production expressions.}
\label{tab:mba}
\begin{tabularx}{\textwidth}{@{}>{\raggedright\arraybackslash}p{0.30\textwidth}X@{}}
\toprule
Tool / corpus & Result \\
\midrule
SiMBA, linear baseline & 18/18 verified reductions \\
GAMBA, linear baseline & 18/18 verified reductions \\
GAMBA, inverse-only nonlinear corpus & 6/6 timeouts or tool failures at 20 s \\
GAMBA, variable-shift nonlinear corpus & 6/6 parser failures \\
SiMBA, inverse-only corpus & 6/6 equivalence-verified reductions (linear simplifier) \\
MBA-Blast & 3/9 verified; 6 parser/tool failures on its grammar subset \\
\bottomrule
\end{tabularx}
\end{table}

The linear forms are fully reducible by both general and linear tools. The
modular-inverse wrapper causes GAMBA to time out; the variable-shift form is
rejected at parse time. SiMBA still reduces the inverse form because it is a
linear simplifier operating on the linear subexpression. MBA-Blast handles a
subset of the shift-free linear cases and rejects the rest. These are measured
cost asymmetries, not a proof of irreducibility.

\subsection{Timing analysis}

The timing evaluation has two layers: wall-clock distinguishability with
dudect-bencher and deterministic instruction/address counts with Cachegrind.

The final dudect run uses an opaque key in the two-key cipher benchmark. Without
that control, the compiler specialized two constant-key paths, and the harness
measured a code-layout difference of about \code{|t| = 19} rather than cipher
data timing. With the opaque key the measured values are:

\begin{center}
\small
\begin{tabular}{@{}lr@{}}
\toprule
Benchmark & \code{max |t|} \\
\midrule
Cipher, two opaque keys & 2.39 \\
Cipher, same key control & 2.05 \\
Normal branchy VM & 204.39 \\
Constant-time straight-line & 1.95 \\
Constant-time branch & 1.67 \\
Constant-time loop & 3.19 \\
Constant-time private memory & 2.73 \\
\bottomrule
\end{tabular}
\end{center}

The four constant-time cases are below the conventional threshold of 5; the
normal branchy profile is far above it. The result confirms the intended
distinction: the default profile is performance-oriented and data-dependent,
while the explicit mode is a fixed-pass oblivious subset.

Cachegrind reports identical counts for the branch-only constant-time leaf across
two inputs: 14,532,847 instruction references, 5,065,486 data references, 4,869
D1 misses, 2,812 LLd misses, 1,339,790 branches, and 110,104 mispredicts. These
counts are deterministic for the tested binary and inputs; they do not cover
wall-clock effects or other microarchitectural state.

\subsection{Performance}

The benchmark protects the same leaf function with Sirus and, separately, with
VMProtect 3.11 Linux demo. Each sample measures 1,000 protected calls; seven
samples are taken and the median is reported. The native and protected checksums
match. Table~\ref{tab:perf} reproduces the measured values.

\begin{table}[htbp]
\centering
\small
\caption{Per-call overhead on the benchmark leaf function (median of seven
samples; 1,000 calls per sample).}
\label{tab:perf}
\begin{tabular}{@{}lrrr@{}}
\toprule
Build & Per call & vs native & vs VMProtect \\
\midrule
Native (Sirus build) & 1.2--1.4 ns & 1$\times$ & 0.0002$\times$ \\
VMProtect 3.11 demo & 8.48 $\mu$s & 4,490$\times$ & 1$\times$ \\
Sirus fast/minimal & 14.5 $\mu$s & 6,100$\times$ & 1.7$\times$ \\
Sirus balanced & 53.9 $\mu$s & 30,000$\times$ & 6.7$\times$ \\
Sirus default & 75.0 $\mu$s & 49,800$\times$ & 9.3$\times$ \\
Sirus hardened & 78.5 $\mu$s & 40,600$\times$ & 9.0$\times$ \\
Sirus constant-time & 234 $\mu$s & 101,500$\times$ & 25.2$\times$ \\
Sirus hardened constant-time & 235 $\mu$s & 111,300$\times$ & 25.5$\times$ \\
\bottomrule
\end{tabular}
\end{table}

The relative-to-native columns use the native baseline of the corresponding
build; the two native measurements differ by about 0.5 ns, so the ratios are
rounded and should be read as engineering estimates, not precise multipliers.

Cold startup is about 0.6 ms for Sirus and about 2.6 s for the packed VMProtect
demo on the same host. Sirus appends about 30 KiB to the benchmark binary
(0.7\%); the VMProtect demo packs the 4.4 MiB input to about 1.17 MiB. The
comparison is deliberately narrow: the VMProtect demo is a real commercial
virtualizer running the same function, but its packing, anti-debug, and demo
checks are not equivalent to Sirus's default profile.

\subsection{Coverage}

Two independent coverage tools are used. Tarpaulin reports 87.07\% line coverage
(7607 of 8737 lines) with the target/bench exclusions. llvm-cov reports 88.25\%
of lines (1786 missed of 15206), 84.09\% of regions (4449 missed of 27962), and
73.53\% of functions (292 missed of 1103). The agreed floor is 85\% lines, and
the proof script enforces it. The uncovered lines are predominantly defensive
error returns, malformed-binary branches, CLI/pipeline I/O failures, and
platform-specific code not executed on Linux.

\subsection{Windows evaluation}

The Windows gate has three jobs:

\begin{enumerate}
  \item \textbf{Wine}: cross-build the protector and fixtures, protect a PE32+
        leaf function, and execute it under Wine in normal, constant-time, and
        hardened modes.
  \item \textbf{Native Windows}: execute the same fixtures on a Windows runner,
        including the ABI matrix, indirect callers, a 32 KiB stack frame, DLL
        export loading, TLS callbacks, and \code{RCDATA} resources.
  \item \textbf{MSVC \code{/guard:cf}}: build with Control Flow Guard, resolve
        function sizes from the PDB and linker map, protect by explicit range,
        and execute the CFG binary with indirect callers.
\end{enumerate}

All three jobs pass on the release tree. The Windows path is a selective leaf
subset, not full Win64 virtualization. External imports, runtime exception
propagation, signing, and kernel anti-debug are out of scope.

\subsection{Threats to validity}

\begin{itemize}
  \item \textbf{Tool scope.} Ghidra, radare2, angr, KLEE, and the MBA tools are
        run with fixed versions and budgets. A different version or a longer
        budget can produce different results. A timeout is not a proof of
        impossibility.
  \item \textbf{Commercial comparison.} VMProtect is a closed product; the
        benchmark uses its Linux demo and a single function. Its security
        profile is not equivalent to Sirus, and the comparison is a performance
        datapoint, not a security ranking.
  \item \textbf{Benchmark function.} The performance leaf is small and
        arithmetic-heavy. Larger functions, memory-heavy loops, or many direct
        calls have different costs.
  \item \textbf{Timing host.} dudect and Cachegrind results are specific to the
        host, compiler, and microarchitecture. The paper reports the measured
        numbers and the exact harness, not a universal constant-time claim.
  \item \textbf{Symbolic resistance.} angr and KLEE time out. The paper does not
        claim that no symbolic path exists.
  \item \textbf{Formal scope.} Kani, Miri, and Creusot cover bounded components
        or models, not the entire product. The decoder is too large for a
        useful CBMC run.
  \item \textbf{Independent audit.} No external audit has been performed. The
        evidence is reproducible, but reproducibility is not the same as
        independent review.
\end{itemize}
